\section{Verification, calibration and errors found}
\label{sec:verification}

A model of this kind is only as good as the evidence that it computes what it
claims. This section records the verification strategy, the calibration against
measured hardware, and, because they are instructive, the errors that the
strategy caught.

\subsection{Verification against the closed forms}

The closed forms of \cref{sec:fold,sec:scaling} exist in the implementation
alongside the numerical solvers and serve as an oracle. The test suite
comprises 66 assertions; those bearing on the theory are:

\begin{itemize}[nosep]
  \item \emph{Momentum theory.} $T = 2\rho A\vind^2$ holds identically to
        machine precision; with $\Cdz\to0$, $\kappa\to1$ and the near-field
        corrections disabled, the computed $\FM\to1$ to within $10^{-5}$,
        confirming \cref{thm:momentum,eq:fm}.
  \item \emph{Coefficients.} $\alpha$ and $\beta$ from \cref{def:coeffs} agree
        with hand computation to four figures.
  \item \emph{Fold.} $m^*$ and $\Mmax$ agree with \cref{eq:closedform} to
        $10^{-12}$ relative; both roots satisfy \cref{eq:closure} to $10^{-8}$;
        $m_-<m^*<m_+$ as \cref{thm:trichotomy} requires; at
        $\Mzero=\Mmax-10^{-9}$ the two roots agree to $10^{-3}$ relative,
        exhibiting the square-root merging of \cref{eq:sqrtlaw}; and for
        $\Mzero>\Mmax$ no solution is returned.
  \item \emph{Iteration.} The naive iteration converges to the analytic $m_-$
        to $10^{-4}$ relative from a distant start, and is classified as
        diverging when $\Mzero>\Mmax$, confirming \cref{thm:iteration}.
  \item \emph{Scaling.} $\Mmax(2t_f) = \Mmax(t_f)/4$ to $10^{-9}$, and the
        numerically computed logarithmic sensitivities reproduce
        \cref{thm:sensitivity} to three decimal places, including the
        non-obvious $S_g=-3.770$ which depends on $\gamma_m/\alpha$.
  \item \emph{Exponent.} The endurance wall is finite for $p=0$ and absent for
        $p=1.4$, confirming \cref{thm:pfold}.
  \item \emph{Dynamics.} With zero thrust the vertical acceleration is exactly
        $-g$; at the hover trim the allocation returns $T=mg$ and
        $\vect\tau=\vect0$ to $10^{-9}$; $\dt E$ equals the negative of the
        computed electrical power identically; rotations satisfy
        $RR\transp=I$ and $\det R=1$ to $10^{-10}$ through the quaternion
        round trip; and the hover linearisation is controllable at full rank.
  \item \emph{Review regressions.} The near-fold root pair of
        \cref{sec:numerics} is found; the quadrotor authority equals the
        balanced differential \cref{eq:authquad} on every axis; the slab
        inertia is assigned to the panel axes; a forward boom produces a
        product of inertia and a trimmed hover with unequal rotor speeds;
        telemetry at $t$ compares the state at $t$ with the references at $t$;
        the governor keeps a reference out of a head that walks toward it;
        the battery loop's reported battery survives its verification flight;
        and the inner-map slope of \cref{prop:inner} lies in $[0,0.9)$ over
        \SIrange{0.8}{50}{\kilo\gram} for the shipped component models.
\end{itemize}

\subsection{Calibration of the acoustic anchor, and a fourteen-decibel error}

\Cref{mod:spl} contains an anchor $C$ which was initially assigned by
judgement, at $C=68$. Calibration against published declarations exposed this
as badly wrong.

Manufacturers declare a sound \emph{power} level $L_{WA}$ under EU regulation
2019/945. Converting to pressure at \SI{1}{\metre} by \cref{eq:eightdb} and
inverting \cref{eq:splmodel} gives $C$ for each aircraft.

\begin{table}[t]
\centering
\caption{Calibration of the acoustic anchor. Hover rotor speeds are
manufacturer figures; tip speed and thrust per rotor follow from geometry and
mass. The four aircraft span a factor of $3.6$ in mass and agree on $C$ to
\SI{1.5}{\decibel}.}
\label{tab:acoustic}
\begin{tabular}{@{}lrrrrrr@{}}
\toprule
Aircraft & $m$ & $T/\text{rotor}$ & $\Vtip$ & $L_{WA}$ & $L_p(\SI{1}{\metre})$ & $C$ \\
\midrule
DJI Mini 3      & \SI{0.249}{\kilo\gram} & \SI{0.61}{\newton} & \SI{71.6}{\metre\per\second} & 79.0 & 71.0 & 82.8 \\
DJI Air 2S      & \SI{0.595}{\kilo\gram} & \SI{1.46}{\newton} & \SI{67.1}{\metre\per\second} & 82.0 & 74.0 & 83.8 \\
DJI Mavic 3     & \SI{0.895}{\kilo\gram} & \SI{2.19}{\newton} & \SI{68.8}{\metre\per\second} & 83.0 & 75.0 & 82.3 \\
DJI Mavic 2 Pro & \SI{0.907}{\kilo\gram} & \SI{2.22}{\newton} & \SI{67.9}{\metre\per\second} & 83.0 & 75.0 & 82.6 \\
\midrule
\multicolumn{6}{r}{mean} & \textbf{82.9} \\
\bottomrule
\end{tabular}
\end{table}

\Cref{tab:acoustic} gives $C = 82.9$ with a spread of \SI{1.5}{\decibel}. The
consistency across a $3.6\times$ mass range is strong evidence that the
functional form \cref{eq:splmodel} is right, and that the original anchor was
\SI{14}{\decibel} optimistic. \Cref{fig:acoustic} shows the four values against
the one first assigned.

\begin{figure}[tbp]
\centering
% Calibration of the acoustic anchor C: each aircraft's declared sound power
% inverted through eq. (splmodel), against the value first assigned.
\begin{tikzpicture}
  \begin{axis}[paperplot, height=5.4cm,
      xmin=0.15, xmax=1.0, ymin=64, ymax=88,
      xlabel={aircraft mass [\si{\kilo\gram}]},
      ylabel={anchor $C$ [\si{\decibel}]},
      legend style={at={(0.98,0.36)}, anchor=east}]
    \fill[cA!10] (axis cs:0.15,82.3) rectangle (axis cs:1.0,83.8);
    \addplot[only marks, mark=*, mark size=2.6pt, cA] coordinates
      {(0.249,82.8) (0.595,83.8) (0.895,82.3) (0.907,82.6)};
    \addlegendentry{calibrated from $L_{WA}$}
    \addplot[cA, dashed, domain=0.15:1.0] {82.9};
    \addlegendentry{mean, $C=82.9$}
    \addplot[cD, line width=1.1pt, domain=0.15:1.0] {68};
    \addlegendentry{initial judgement, $C=68$}
    \node[lbl, anchor=south] at (axis cs:0.249,83.0) {Mini 3};
    \node[lbl, anchor=south] at (axis cs:0.595,84.0) {Air 2S};
    \node[lbl, anchor=north east] at (axis cs:0.88,82.1) {Mavic 3};
    \node[lbl, anchor=south west] at (axis cs:0.91,82.8) {Mavic 2};
    \draw[dim, cD] (axis cs:0.42,68) -- (axis cs:0.42,82.9)
      node[midway, right, font=\footnotesize, cD, align=left]
      {the anchor was\\ this far optimistic};
  \end{axis}
\end{tikzpicture}

\caption{The acoustic anchor $C$ recovered from the declared sound power of
four aircraft through \cref{eq:eightdb,eq:splmodel}, against aircraft mass.
The four values lie in the narrow band shaded across a factor of $3.6$ in
mass, about the mean $C=82.9$. The value first assigned by judgement, $C=68$,
lies below every one of them by the length of the arrow, an error larger than
the effect of any design choice in \cref{sec:results}.}
\label{fig:acoustic}
\end{figure}

\begin{finding}
Fourteen decibels is not a refinement. It is the difference between a machine
one can think beside and one that must leave the room, and it exceeds the
entire range spanned by every design decision available in the study. The error
had two components: an uncalibrated anchor, and the eight-decibel confusion of
\cref{prop:lwlp} between a declared sound power and a pressure at one metre. It
was found only because the model was made to reproduce measured hardware.
\end{finding}

The test suite now asserts that \cref{eq:splmodel} reproduces the declared
levels of three of these aircraft to within \SI{1.5}{\decibel}.

\subsection{Calibration of the motor model}

$k_\tau$ in \cref{eq:motormass} was likewise assigned by judgement at
\SI{3.0}{\newton\metre\per\kilo\gram} and found to be pessimistic by roughly a
factor of two.

\begin{table}[t]
\centering
\caption{Motor mass model against catalogue peak torque. The model is
conservative at the large end, which is the safe direction for a design study;
the designs here use motors of \SIrange{50}{60}{\gram} where the error is
small.}
\label{tab:motor}
\begin{tabular}{@{}lrrrr@{}}
\toprule
Motor & Actual mass & Peak torque & Model mass & Error \\
\midrule
T-Motor MN2806  & \SI{66}{\gram}  & \SI{0.30}{\newton\metre} & \SI{63}{\gram}  & $-5\%$ \\
T-Motor MN3110  & \SI{96}{\gram}  & \SI{0.63}{\newton\metre} & \SI{111}{\gram} & $+16\%$ \\
T-Motor MN505-S & \SI{190}{\gram} & \SI{2.03}{\newton\metre} & \SI{273}{\gram} & $+44\%$ \\
\bottomrule
\end{tabular}
\end{table}

With $k_\tau = \SI{5.5}{\newton\metre\per\kilo\gram}$ at
$m_{\mathrm{ref}}=\SI{0.1}{\kilo\gram}$ and $e=0.30$, \cref{tab:motor} results.
The recalibration mattered structurally, not just numerically: with the
pessimistic constant the motor mass penalty overwhelmed the induced-power
benefit of a larger rotor, and the optimum diameter sat on the lower bound of
the search box. After recalibration a genuine interior optimum appears, which
is the physically expected behaviour, and the test suite now asserts its
existence. \Cref{fig:motorcal} shows the calibrated law against the catalogue.

\begin{figure}[tbp]
\centering
% The motor mass model m = (Q m_ref^e / k_tau)^{1/(1+e)} with k_tau = 5.5
% N m/kg, m_ref = 0.1 kg, e = 0.30, against catalogue motors; the connectors
% are the model errors of tab:motor.
\providecommand{\fdQuadRoll}{3.66}
\providecommand{\fdQuadPitch}{3.66}
\providecommand{\fdQuadOld}{0.91}
\providecommand{\fdHexRoll}{6.13}
\providecommand{\fdHexPitch}{7.07}
\providecommand{\fdHexGross}{2.565}
\providecommand{\fdQuadGross}{1.724}
\providecommand{\fdKeepRadius}{1.180}
\providecommand{\fdKeepSlice}{1.174}
\providecommand{\fdEnvRadius}{1.080}
\providecommand{\fdTurnT}{5.65}
\providecommand{\fdTurnHeadX}{-1.440}
\providecommand{\fdTurnHeadY}{0.000}
\providecommand{\fdTurnPeriod}{7.54}
\providecommand{\fdStandoff}{1.20}
\providecommand{\fdRootLo}{4.5678}
\providecommand{\fdRootHi}{4.6173}
\providecommand{\fdPeakM}{4.5925}
\providecommand{\fdPeakF}{0.03}
\providecommand{\fdGridLo}{4.460}
\providecommand{\fdGridHi}{6.021}
\providecommand{\fdGridLoF}{-0.001}
\providecommand{\fdGridHiF}{-0.083}
\providecommand{\fdIslandPct}{1.1}
\providecommand{\fdPlainRate}{0.263}
\providecommand{\fdConvWindow}{8}
\providecommand{\fdHlMotorRe}{-0.0889}
\providecommand{\fdHlMotorIm}{0.0000}
\providecommand{\fdHlAttRe}{-0.0240}
\providecommand{\fdHlAttIm}{0.0000}
\providecommand{\fdHlPosRe}{-0.0082}
\providecommand{\fdHlPosIm}{0.0004}
\providecommand{\fdHlGovRe}{-0.0160}
\providecommand{\fdHlGovIm}{0.0000}
\providecommand{\fdWpos}{2.04}
\providecommand{\fdWgov}{4.00}
\providecommand{\fdWatt}{6.43}
\providecommand{\fdWmot}{22.2}
\providecommand{\fdWarm}{206}
\providecommand{\fdWnyq}{785}
\providecommand{\fdRatioPA}{3.2}
\providecommand{\fdRatioAM}{3.5}
\providecommand{\fdKR}{41.4}
\providecommand{\fdKp}{4.17}
\providecommand{\fdTauMms}{45}
\providecommand{\fdH}{4}
\providecommand{\fdMotorQ}{0.289}
\providecommand{\fdMotorM}{0.0609}

\begin{tikzpicture}
  \begin{axis}[paperplot, height=5.6cm, xmode=log, ymode=log,
      xmin=0.1, xmax=3.0, ymin=20, ymax=400,
      xlabel={peak torque $Q$ [\si{\newton\metre}]},
      ylabel={motor mass [\si{\gram}]},
      xtick={0.1,0.2,0.5,1,2,3}, xticklabels={0.1,0.2,0.5,1,2,3},
      ytick={20,50,100,200,400}, yticklabels={20,50,100,200,400},
      legend pos=north west]
    \addplot[cA, domain=0.1:3.0, samples=60]
      {1000*(x*0.1^0.3/5.5)^(1/1.3)};
    \addlegendentry{model \labelcref{eq:motormass}, $k_\tau=\SI{5.5}{\newton\metre\per\kilo\gram}$, $e=0.30$}
    \addplot[only marks, mark=square*, mark size=2.4pt, cB] coordinates
      {(0.30,66) (0.63,96) (2.03,190)};
    \addlegendentry{catalogue motors, T-Motor}
    \addplot[only marks, mark=*, mark size=2.8pt, cC] coordinates
      {(\fdMotorQ,1000*\fdMotorM)};
    \addlegendentry{the design's motor}
    \draw[cD, line width=0.8pt] (axis cs:0.30,66) -- (axis cs:0.30,63);
    \draw[cD, line width=0.8pt] (axis cs:0.63,96) -- (axis cs:0.63,111);
    \draw[cD, line width=0.8pt] (axis cs:2.03,190) -- (axis cs:2.03,273);
    \node[lbl, anchor=north west] at (axis cs:0.31,64) {MN2806, $-5\%$};
    \node[lbl, anchor=north west] at (axis cs:0.65,94) {MN3110, $+16\%$};
    \node[lbl, anchor=north] at (axis cs:2.03,178) {MN505-S, $+44\%$};
  \end{axis}
\end{tikzpicture}

\caption{The calibrated motor-mass model \cref{eq:motormass} against catalogue
motors, on logarithmic axes on which the model is a straight line of slope
$1/(1+e)=0.77$. The red connectors are the model errors of \cref{tab:motor}:
small at the size of motor this design uses (the green point, the design's
\SI{61}{\gram} motor at its torque for maximum thrust) and conservative at the
large end.}
\label{fig:motorcal}
\end{figure}

\subsection{Errors found by the model itself}

Three defects were found not by inspection but because a computed quantity was
inconsistent with a hand calculation.

\paragraph{The tip-speed design rule.} Holding tip speed fixed while growing
the rotor made power \emph{increase} with diameter, contradicting
\cref{obs:dl}. The cause was that at fixed tip speed the profile power
\cref{eq:profile} scales with disk area while the induced power falls only as
$A^{-1/2}$. The resolution is \cref{cor:designrule}: a rotor is designed at
constant blade loading, not constant tip speed, and with $\Vtip$ derived from
\cref{eq:ctsigma} the expected monotonicity is recovered.

\paragraph{Safety measured to the wrong point.} The simulator's minimum-distance
check measured from the vehicle centre of mass to the user. By \cref{eq:reach}
the blades extend $\mathcal{R}$ from that point, and for these designs
$\mathcal{R}=\SI{0.63}{\metre}$, so the check was optimistic by the entire reach
of the machine. Recomputing to the nearest blade tip showed the nominal
configuration violating its own safety distance by a factor of three. This is
discussed in \cref{sec:results}.

\paragraph{A startup transient masquerading as a failure.} The simulation began
with the vehicle at rest at the reference position while the human was already
at full walking speed. On a curved path the display sits directly in the
person's path, so the separation collapsed from \SI{1.100}{\metre} to
\SI{0.603}{\metre} within \SI{0.62}{\second} before recovering. Initialising the
vehicle and the governor with the reference velocity, as noted in
\cref{sec:numerics}, removed it; the minimum envelope gap then agrees with the
static tip gap of \cref{eq:clearance} to within a centimetre, the difference
being the conservative sphere against the actual blade tips. The same mistake
reappeared one level up: the vehicle was initialised level while the
reference on a curved path already needed a tilt and a body rate, and the
attitude loop spent the first second catching up. That transient set the
whole-window peak tracking error, which is the quantity the keep-out margin is
checked against. The initial attitude, body rate and thrust are now the trim
for the reference at $t=0$.

\begin{remark}
The common feature of all of these is that they were caught by comparison
against an independently computed quantity, not by reading the code. This is
the argument for keeping the analytic layer alive in the implementation even
after it has been superseded.
\end{remark}

\subsection{Errors found by independent review}

An independent review of the engine and of an earlier draft of this paper
found five defects that the tests above had not, each now covered by a
regression test.

\begin{enumerate}[nosep]
  \item The component closure bracketed the first sign change of
        $\mathcal{F}$ on a geometric grid of ratio $1.35$ and called a grid
        with no positive sample infeasible. Near the fold the positive stretch
        is narrower than the grid, and a design with two roots was reported as
        having none. The paper had also described that failure as
        infeasibility. \Cref{sec:numerics} gives the replacement.
  \item The torque authority divided the balanced differential by $N$, a
        factor of four on every axis, which propagated into the gain rule and
        into every layout comparison. \Cref{eq:authority} is now a linear
        program and \cref{app:proofs} the closed form it reduces to.
  \item Telemetry compared the state advanced to $t+\dd t$ with references
        evaluated at $t$, labelled $t$.
  \item The display slab's inertia was assigned to the wrong body axes, and the
        offsets were not moved to the centre of mass, so a boom produced no
        product of inertia. \Cref{sec:components} now builds one tensor from
        one geometry.
  \item After the review, and found by the tracking statistics it demanded:
        the attitude loop used a zero desired angular rate, so on a curved
        path it lagged the rotating thrust direction by a standing error. The
        rate feedforward of \cref{eq:omegad} removed a \SI{0.27}{\metre}
        tracking excursion and a rotor pass at \SI{0.34}{\metre} from the eye.
\end{enumerate}
